%! Right Angles Revisited --- Setting Up Orthogonality — Unit 4, Lesson 4.0 — Exit Ticket
\documentclass[10pt]{article}
\usepackage{linalg-article}
\usepackage{linalg-boxes}
\newcommand{\TallMath}[1]{$\displaystyle #1\rule[-1.4em]{0pt}{3.2em}$}
\newcommand{\vv}{\mathbf{v}}
\newcommand{\ww}{\mathbf{w}}
\newcommand{\zero}{\mathbf{0}}

\begin{document}

\pageheader{Unit 4, Lesson 4.0}{Exit Ticket}
\namedateperiod

\vspace{0.15in}

No notes. Show your reasoning.

\begin{enumerate}[leftmargin=1.8em, itemsep=16pt]
  \item \textbf{Dot product \& right angle.} Compute
        $\begin{bmatrix} 4 \\ 1 \end{bmatrix}\cdot\begin{bmatrix} 2 \\ -8 \end{bmatrix}$.
        Are the two vectors perpendicular? How do you know?
        \vspace{1.1cm}
  \item \textbf{Length \& unit vector.} For $\begin{bmatrix} 5 \\ 12 \end{bmatrix}$, find the length
        $\|\cdot\|=\sqrt{\;\cdot\;}$ and the unit vector in the same direction.
        \vspace{1.3cm}
  \item \textbf{Solve \& justify.} Find the value of $x$ that makes
        $\begin{bmatrix} 3 \\ x \end{bmatrix}$ perpendicular to $\begin{bmatrix} 6 \\ -2 \end{bmatrix}$.
        Then explain in one sentence what $\vv\cdot\ww=0$ tells you about the \emph{angle} between two
        vectors.
        \writelines{2}
\end{enumerate}

\end{document}
